\section*{Chapter 8 --- SPI and the transaction protocol}
\addcontentsline{toc}{section}{Chapter 8 --- SPI and the transaction protocol}

\solution{ex:8:1}{Decode}
The command byte is \code{W} in bit 7, zeroes in bits 6--4, and the register index in bits 3--0.
\ans{a} \code{0x81} = write bit set, index 1: a \textbf{write to \code{TX_ID}} with the value
\code{0x000007FF} --- eleven bits of ones, that is, the largest value the register can hold.
\ans{b} \code{0x85} = write, index 5: a \textbf{write to \code{TX_SEND}} with \code{0x1}, which
\emph{triggers a transmission}. The register stores nothing.
\ans{c} \code{0x8B} = write, index 11: a \textbf{write to \code{ERROR_FLAGS}} with \code{0x0},
which clears the error bit.
\ans{d} \code{0xC2} is \code{1100 0010}: the write bit set, but \emph{also} bit 6, which is
reserved and has to be zero. The command is invalid --- no write is committed, no read is latched
--- and the slave consumes all five bytes anyway before returning to idle.

The two unusual ones are \textbf{c} and \textbf{d}. In c you clear by writing \emph{zero}, which
is the opposite rule to \code{TX_SEND} and \code{RX_ACK}, which require bit 0 to be set. In d the
command is invalid and the slave does \emph{not} report it: there is no error channel at this
layer, so a driver sending junk gets silent nonsense back.

\solution{ex:8:2}{Encode}
\ans{a} A read has bit 7 low, so the command byte is just the index; \code{RX_DLC} has index 7 and
the four data bytes are dummy bytes: \code{07 00 00 00 00}
\ans{b} \code{82 00 00 00 04}
\ans{c} \code{85 00 00 00 01}
\ans{d} \code{8A 00 00 00 01}
\ans{e} \code{8B 00 00 00 00}

The difference between d and e is the rule for what counts as a valid write. \code{RX_ACK}
requires \textbf{bit 0 to be set} --- a zero does nothing. \code{ERROR_FLAGS} requires \textbf{the
whole word to be zero} --- \code{0x1} does nothing. Two registers next to each other in the same
map, with opposite rules, and confusing them is one of the most common bugs in the project.

\solution{ex:8:3}{What comes back}
\ans{a} None. During a write, everything the slave drives on \code{MISO} is meaningless.
\ans{b} Four: the last four. The byte clocked in \emph{during the command byte} is leftovers from
last time, since the slave does not load its response shifter until the end of the command byte.
\ans{c} The return value is shifted by one byte: the first real data byte ends up in bits 23--16
instead of 31--24, and the last data byte falls off entirely. While debugging it looks as though
every register value is $256$ times too small, or --- for \code{STATUS} --- as though the bit you
are looking for is never set. That is the same symptom as when \code{SPI0.DATA} is not read in
\code{transfer()}, which makes the two easy to confuse.

\solution{ex:8:4}{\texorpdfstring{\code{SS}}{SS} and aborts}
\ans{a} The transaction is abandoned entirely, without side effects: \code{TX_ID} is not changed,
no trigger fires, and the slave returns to idle ready for the next falling \code{SS} edge. That is
what makes the slave self-healing after a glitch.
\ans{b} The slave takes the first five bytes as a transaction and \textbf{discards the following
five}. A transaction only begins on the first byte after \code{SS} has fallen, so the sixth byte
is never interpreted as a new command byte.

The dangerous part is that it fails \emph{silently}. At bring-up it looks as though every other
register write disappears: the first arrives, the second does not, and nothing reports an error.
You suspect the wiring, the clock frequency and the driver's bit packing --- long before you
suspect \code{SS}.
\ans{c} \code{0x91} is \code{1001 0001}: the write bit set, index 1, but bit 4 is set and
reserved. The command is invalid and is ignored entirely. It should be prevented already when the
command byte is built: \code{CmdWrite | index} with an index in the range 0--12 can never set bits
6--4, which is why the index is a named constant and not a number from a calculation.

\solution{ex:8:5}{Time budget}
\ans{a} Six transactions --- \code{STATUS}, \code{RX_ID}, \code{RX_DLC}, two data registers and
\code{RX_ACK} --- that is, around $6 \times 45 = 270$~µs.
\ans{b} A maximal frame is around 130 bits, that is, about \textbf{130~µs} at 1~Mbit/s.
\ans{c} They overwrite the previous one. The controller holds exactly one received frame, and the
newest wins.

That is \textbf{not} your bug: it is the hardware's documented limitation, and it is deliberate.
You can reduce the risk by polling often and not doing anything slow between the calls, but you
cannot get past the fact that a \code{receive()} takes twice as long as the frame it fetches. At
sufficient traffic intensity frames are lost whatever the driver does.
